\BourbakiExercisesSection

\begin{enumerate}
\def\labelenumi{\arabic{enumi}.}
\tightlist
\item
  In a finite group G, show that the number of conjugates (no. 4) of an element \(a \in G\) is equal to the index of the normalizer of a and is therefore a divisor of the order of G.
\end{enumerate}

¶ 2. *If G is a finite group of order n, the number of automorphisms of G is \(\leqslant n^{\log n/\log 2}\) (show that there exists a generating set \(\{a_{1}, a_{2}, \ldots, a_{m}\}\) of G such that \(a_{i}\) does not belong to the subgroup generated by \(a_{1}, a_{2}, \ldots, a_{i-1}\) for \(2 \leqslant i \leqslant m\) ; deduce that \(2^{m} \leqslant n\) and that the number of automorphisms of G is \(\leqslant n^{m})\cdot*\)

\begin{enumerate}
\def\labelenumi{\arabic{enumi}.}
\setcounter{enumi}{2}
\tightlist
\item
  Let \(\Gamma\) be the group of automorphisms of a group \(G\) and \(\Delta\) the group of inner automorphisms of \(G\) ; show that \(\Delta\) is a normal subgroup of \(\Gamma\) . For an automorphism \(\sigma\) of \(G\) to be permutable with all the inner automorphisms of \(G\) , it is necessary and sufficient that, for all \(x \in G\) , \(x^{-1}\sigma(x)\) belong to the centre of \(G\) ; deduce that, if the centre of \(G\) is reduced to the identity element, so is the centralizer of \(\Delta\) in \(\Gamma\) .
\end{enumerate}

¶ 4. (a) Let G be a non-commutative simple group, Γ the automorphism group of G and Δ the group of inner automorphisms of G (isomorphic to G). If s is an automorphism of the group Γ, show that \(s(\Delta) = \Delta\) (using Exercise 3 above and § 4, no. 9, Proposition 15, note that the intersection \(\Delta \cap s(\Delta)\) cannot consist only of the identity element of Γ).

\begin{enumerate}
\def\labelenumi{(\alph{enumi})}
\setcounter{enumi}{1}
\item
  Show that the only automorphism of \(\Gamma\) leaving invariant each of the elements of \(\Delta\) is the identity element (use the fact that this automorphism leaves invariant \(\alpha\) and \(\sigma \alpha \sigma^{-1}\) for all \(\alpha \in \Delta\) and \(\sigma \in \Gamma\) and use Exercise 3).
\item
  Let \(s\) be an automorphism of \(\Gamma\) , \(\phi\) the isomorphism \(x \mapsto \operatorname{Int}(x)\) of G onto \(\Delta\) , \(\psi\) the inverse isomorphism and \(\sigma\) the automorphism \(\psi \circ s \circ \phi\) of G; show that the automorphism \(\xi \mapsto \sigma^{-1}s(\xi)\sigma\) of \(\Gamma\) is the identity automorphism (use (b), noting that, for all \(x \in G\) , \(s(\operatorname{Int}(x)) = \operatorname{Int}(\sigma(x)))\) . Deduce that every automorphism of \(\Gamma\) is an inner automorphism.
\end{enumerate}

5. Let G be a group.

\begin{enumerate}
\def\labelenumi{(\alph{enumi})}
\item
  If H is a subgroup of G of finite index n, show that the intersection N of the conjugates of H is of index a divisor of \(n!\) (note that G/N is isomorphic to a subgroup of \(S_{n}\) ).
\item
  G is called residually finite, if the intersection of its subgroups of finite index is \(\{e\}\) . Show that this is equivalent to saying that G is isomorphic to a subgroup of a product of finite groups.
\end{enumerate}
% semantic-fragments: legacy-input-seam
\relax{}
\begin{enumerate}
\def\labelenumi{(\alph{enumi})}
\setcounter{enumi}{2}
\item
  Suppose that G can be generated by a finite set. Show that, for every integer n, the set \(P_{n}\) of subgroups of G of index n is finite.
\item
  Under the hypothesis of (c), consider an endomorphism \(f\) of \(G\) which is surjective. Show that, for every integer \(n\) , the mapping \(H \mapsto f^{-1}(H)\) is a bijection of \(P_n\) onto itself (observe that it is an injection). Deduce that the kernel of \(f\) is contained in every subgroup of \(G\) of finite index. In particular, if \(G\) is residually finite, \(f\) is bijective.
\end{enumerate}

\begin{enumerate}
\def\labelenumi{\arabic{enumi}.}
\setcounter{enumi}{5}
\tightlist
\item
  Let H be a subgroup of finite index in a group G. Suppose that G is the union of the conjugates of H. Show that H = G. (Reduce it to the case where G is finite by means of Exercise 5. In the latter case, note that:
\end{enumerate}

\[
\operatorname{Card} \left(\bigcup_ {x \in G} x H x ^ {- 1}\right) \leqslant 1 + (\operatorname{Card} (H) - 1) \operatorname{Card} (G / H).)
\]

\begin{enumerate}
\def\labelenumi{\arabic{enumi}.}
\setcounter{enumi}{6}
\tightlist
\item
  Let \(p\) be a prime number.
\end{enumerate}

\begin{enumerate}
\def\labelenumi{(\alph{enumi})}
\item
  Show that \(n^p \equiv n \pmod{p}\) for all \(n \in \mathbf{Z}\) . (Reduce it to the case where \(n\) is positive and argue by induction on \(n\) using the binomial formula.)
\item
  Let \(x, y\) be two elements of a group \(G\) . Suppose that
\end{enumerate}

\[
y x y ^ {- 1} = x ^ {n}, \quad \text { with } n \in \mathbf {Z}, \quad \text { and   that } \quad x ^ {p} = 1.
\]

Show that \(y^{p}xy^{-p} = x^{n}\) ; deduce that \(y^{p - 1}\) commutes with \(x\) .

\begin{enumerate}
\def\labelenumi{(\alph{enumi})}
\setcounter{enumi}{2}
\tightlist
\item
  Let G be a group all of whose elements \(\neq1\) are of order p and are conjugate to one another. Show that \(\operatorname{Card}(G)\) is equal to 1 or 2. (Use (b) to prove that G is commutative.)
\end{enumerate}

\begin{enumerate}
\def\labelenumi{\arabic{enumi}.}
\setcounter{enumi}{7}
\tightlist
\item
  Let A and B be subgroups of a finite group G, \(N_{A}\) and \(N_{B}\) the normalizers of A and B, \(v_{A}\) and \(v_{B}\) the indices of \(N_{A}\) and \(N_{B}\) in G, \(r_{A}\) the number of conjugates of A which contain B and \(r_{B}\) the number of conjugates of B which contain A; show that \(v_{A}r_{B} = v_{B}r_{A}\) .
\end{enumerate}

¶ 9. Let G be a group of permutations of a finite set E; for all \(s \in G\) , let \(\chi(s)\) denote the cardinal of the set of fixed points of \(s\) .

\begin{enumerate}
\def\labelenumi{(\alph{enumi})}
\item
  Show that \(\sum_{s\in \mathbf{G}}\chi (s) = \mathrm{N}t\) , where \(\mathbf{N} = \operatorname {Card}(\mathbf{G})\) and \(t\) is the number of orbits of \(\mathbf{G}\) in \(\mathbf{E}\) (evaluate in two ways the number of ordered pairs \((s,x)\in \mathbf{G}\times \mathbf{E}\) such that \(s(x) = x\) ).
\item
  Suppose that, for all \(a \in \mathbf{E}\) , the stabilizer \(\mathrm{H}_a\) of \(a\) is not reduced to \(e\) . Show that, if \(\chi(s)\) is independent of \(s\) for \(s \neq e\) and is equal to an integer \(k\) , then
\end{enumerate}

\[
k \leqslant t \leqslant 2 k.
\]

(Note that \(\operatorname{Card}(\mathbf{E}) \leqslant \mathbf{N}k\) .)

In the particular case where k = 2, find the possible orders of the groups \(H_{a}\) ; show that, if t = 3, the order of \(H_{a}\) cannot be \(\geqslant 3\) for elements of two of the three orbits unless N has one of the values 12, 24 or 60.

\begin{enumerate}
\def\labelenumi{(\alph{enumi})}
\setcounter{enumi}{2}
\tightlist
\item
  Suppose that G operates transitively on E and let H denote the stabilizer of an element of E; show that \(\sum_{s\in G}\chi(s)^{2}\) is equal to N. \(t_{H}\) , where \(t_{H}\) is the number of orbits of H in E. (For every \((s,u)\) , write \(\chi(s,u)=\chi(s)\) . Evaluate in two ways the sum \(\sum_{(s,u)\in R}\chi(s,u)\) , where R is the set of ordered pairs \((s,u)\) such that \(usu^{-1}\in H.\) )
\end{enumerate}

\begin{enumerate}
\def\labelenumi{\arabic{enumi}.}
\setcounter{enumi}{9}
\tightlist
\item
  \begin{enumerate}
  \def\labelenumii{(\alph{enumii})}
  \tightlist
  \item
    Show that the elements \(\tau_{12}\tau_{34},\tau_{13}\tau_{24},\tau_{14}\tau_{23}\) of the alternating group \(\mathfrak{A}_4\) form with the identity element a commutative subgroup H of \(\mathfrak{A}_4\) . Show that H is normal in \(\mathfrak{A}_4\) and in \(\mathfrak{S}_4\) and that \(\mathfrak{A}_4 / \mathrm{H}\) is cyclic of order 3.
  \end{enumerate}
\end{enumerate}

\begin{enumerate}
\def\labelenumi{(\alph{enumi})}
\setcounter{enumi}{1}
\item
  Show that the centralizer of H in \(S_{4}\) is equal to H. Deduce that the mapping \(s \mapsto (h \mapsto shs^{-1})\) of \(S_{4}\) into \(\operatorname{Aut}(H)\) defines when passing to the quotient an isomorphism of \(S_{4}/H\) onto \(\operatorname{Aut}(H)\) and that the latter group is isomorphic to \(S_{3}\) .
\item
  Let \(\mathbf{K}\) be a subgroup of \(\mathbf{H}\) of order 2. Show that \(\mathbf{K}\) is not normal in \(\mathfrak{A}_4\) .
\end{enumerate}

\begin{enumerate}
\def\labelenumi{\arabic{enumi}.}
\setcounter{enumi}{10}
\tightlist
\item
  Let E be a finite set, \(\zeta \in \mathfrak{S}_{E}\) a cycle of support E and \(\tau\) a transposition. Show that \(\mathfrak{S}_{E}\) is generated by \(\zeta\) and \(\tau\) .
\end{enumerate}

¶ 12. (a) Show that every permutation \(\sigma \in \mathfrak{A}_n\) is a product of cycles of order 3 (which are not in general component cycles of \(\sigma\) ). (Prove it for a product of two transpositions and use § 5, no. 7, Proposition 9.)

\begin{enumerate}
\def\labelenumi{(\alph{enumi})}
\setcounter{enumi}{1}
\item
  If \(a_1, \ldots, a_p\) are \(p\) distinct elements of \([1, n]\) , let \((a_1 a_2 \ldots a_p)\) denote the cycle of order \(p\) whose support is \(\{a_1, a_2, \ldots, a_p\}\) and which maps \(a_i\) to \(a_{i+1}\) for \(1 \leqslant i \leqslant p-1\) and \(a_p\) to \(a_1\) . Show that \(\mathfrak{A}_n\) is generated by the \(n-2\) cycles (1 2 3), (1 2 4), \ldots, (1 2 n). (Use (a).)
\item
  Deduce that, if \(n\) is odd, \(\mathfrak{A}_n\) is generated by (1 2 3) and (1 2 \ldots{} \(n\) ) and, if \(n\) is even, by (1 2 3) and (2 3 \ldots{} \(n\) ).
\item
  Show that, if a normal subgroup of \(\mathfrak{A}_n\) contains a cycle of order 3, it is identical with \(\mathfrak{A}_n\) (prove that such a subgroup contains all the cycles (12 \(k\) ), \(3 \leqslant k \leqslant n\) ).
\end{enumerate}

¶ 13. Let G be a transitive group of permutations of a set X and H the stabilizer of an element \(x \in X\) . Show the equivalence of the following properties:

\begin{enumerate}
\def\labelenumi{(\alph{enumi})}
\item
  Every subgroup \(\mathbf{H}'\) of \(\mathbf{G}\) containing \(\mathbf{H}\) is equal to \(\mathbf{H}\) or \(\mathbf{G}\) .
\item
  Every subset Y of X such that, for all \(g \in G\) , gY is either contained in Y or disjoint from Y, is equal to X or consists of a single element.
\end{enumerate}

(If \(\mathbf{H}'\) satisfies (a) and \(\mathbf{Y} = \mathbf{H}'x\) , then \(g\mathbf{Y} = \mathbf{Y}\) for all \(g \in \mathbf{H}'\) and \(g\mathbf{Y} \cap \mathbf{Y} = \varnothing\) for all \(g \notin \mathbf{H}'\) . Conversely, if \(\mathbf{Y}\) enjoys the property in (b), the set \(\mathbf{H}'\) of \(g \in \mathbf{G}\) such that \(g\mathbf{Y} = \mathbf{Y}\) is a subgroup of \(\mathbf{G}\) containing \(\mathbf{H}\) .)

A transitive permutation group G satisfying (a) and (b) and not consisting only of the identity element is called primitive.

\begin{enumerate}
\def\labelenumi{\arabic{enumi}.}
\setcounter{enumi}{13}
\tightlist
\item
  A group of permutations \(\Gamma\) of a set \(E\) is called \(r\) -ply transitive if, for any two sequences \((a_1, a_2, \ldots, a_r)\) , \((b_1, b_2, \ldots, b_r)\) of \(r\) distinct elements of \(E\) , there exists a permutation \(\sigma \in \Gamma\) such that \(\sigma(a_i) = b_i\) for \(1 \leqslant i \leqslant r\) , this property not holding for at least one ordered pair of sequences of \(r + 1\) distinct elements of \(E\) .
\end{enumerate}

\begin{enumerate}
\def\labelenumi{(\alph{enumi})}
\item
  Show that an \(r\) -ply transitive group is primitive if \(r > 1\) .
\item
  The order of a permutation group \(\Gamma\) operating on \(n\) objects which is \(r\) -ply transitive is of the form \(n(n - 1)\ldots (n - r + 1)d\) , where \(d\) is a divisor of \((n - r)!\) (consider a subgroup of the permutations of \(\Gamma\) leaving invariant \(r\) elements and calculate its index).
\end{enumerate}

¶ 15. Let \(\Gamma\) be an \(r\) -ply transitive group of permutations of a set \(\mathbf{E}\) with \(n\) elements; for a permutation \(\sigma \in \Gamma\) distinct from the identity permutation, let \(n - s\) be the number of elements of \(\mathbf{E}\) invariant under \(\sigma\) . If \(s > r\) , show that there exists a permutation \(\tau \in \Gamma\) such that \(\sigma^{-1}\tau \sigma \tau^{-1}\) is distinct from the identity permutation and the number of elements of \(\mathbf{E}\) it leaves invariant is \(\geqslant n - 2(s - r + 1)\) (use the decomposition of \(\sigma\) into its component cycles). If \(s = r\) , show similarly that there exists \(\tau \in \Gamma\) such that \(\sigma^{-1}\tau \sigma \tau^{-1}\) is a cycle of order 3. Deduce that, if \(r \geqslant 3\) and \(\Gamma\) does not contain the alternating group \(\mathfrak{A}_n\) , then \(s \geqslant 2r - 2\) for every permutation in \(\Gamma\) (use Exercise 10). Conclude finally that, if \(\Gamma\) is not identical with \(\mathfrak{A}_n\) or \(\mathfrak{S}_n\) , then \(r \leqslant n/3 + 1\) .

¶ 16. (a) Show that the alternating group \(\mathfrak{A}_n\) is \((n - 2)\) -ply transitive.

\begin{enumerate}
\def\labelenumi{(\alph{enumi})}
\setcounter{enumi}{1}
\item
  Show that \(\mathfrak{A}_n\) is a non-commutative simple group for \(n \geqslant 5\) . (Use (a), the method of Exercise 15 and Exercise 12(d) to prove that \(\mathfrak{A}_n\) is simple if \(n > 6\) ; examine analogously the cases \(n = 5\) and \(n = 6\) .)
\item
  Show that, for \(n \geqslant 5\) , the only normal subgroups of \(\mathfrak{S}_n\) are \(\{e\}\) , \(\mathfrak{A}_n\) and \(\mathfrak{S}_n\) .
\end{enumerate}

\begin{enumerate}
\def\labelenumi{\arabic{enumi}.}
\setcounter{enumi}{16}
\item
  Let G be a transitive group of permutations of a set X; suppose G is primitive (Exercise 13). Show that every normal subgroup of G distinct from \(\{e\}\) is transitive.
\item
  Let \(\Gamma\) be a group of permutations of a set E. Let A and B be two subsets of E, complements one of the other and stable under \(\Gamma\) . Let \(\Gamma_{A}\) and \(\Gamma_{B}\) denote the groups consisting of the restrictions of the permutations in \(\Gamma\) to A and B respectively and \(\Delta_{A}\) and \(\Delta_{B}\) the subgroups of \(\Gamma\) which leave invariant respectively every element of A and every element of B. Show that \(\Delta_{A}\) and \(\Delta_{B}\) are normal subgroups of \(\Gamma\) and that \(\Gamma_{A}\) is isomorphic to \(\Gamma/\Delta_{A}\) and \(\Gamma_{B}\) to \(\Gamma/\Delta_{B}\) ; if \(\Delta_{AB}\) (resp. \(\Delta_{BA}\) ) is the group consisting of the restrictions of the permutations in \(\Delta_{A}\) (resp. \(\Delta_{B}\) ) to B (resp. A), show that the quotient groups \(\Gamma_{A}/\Delta_{BA}\) , \(\Gamma_{B}/\Delta_{AB}\) and \(\Gamma/(\Delta_{A}\Delta_{B})\) are isomorphic (use Exercise 8).
\item
  Let G be a group, H a subgroup of G and G/H the set of left cosets (mod. H) in G. Let \(r: \mathrm{G} / \mathrm{H} \to \mathrm{G}\) be a section associated with the canonical projection \(\mathrm{G} \to \mathrm{G} / \mathrm{H}\) . Then for all \(\mathbf{X} \in \mathrm{G} / \mathrm{H}\) , \(\mathbf{X} = r(\mathbf{X})\) . H. An internal law of composition is defined on G/H by setting \(\mathbf{X} \top \mathbf{Y} = r(\mathbf{X})r(\mathbf{Y})\) . H.
\end{enumerate}

\begin{enumerate}
\def\labelenumi{(\alph{enumi})}
\item
  Show that \(X \top H = X\) for all X and that under the law \(\top\) every left translation is bijective. If \(G'\) is the subgroup of G generated by the set of elements \(r(X)\) and \(H' = H \cap G'\) , the internal law defined analogously on \(G'/H'\) by the mapping r determines on this set a structure isomorphic to that determined on G/H by the law \(\top\) .
\item
  For the law \(\top\) to be associative, it is necessary and sufficient that \(\mathbf{H}'\) be a normal subgroup of \(\mathbf{G}'\) in which case the structure determined by \(\top\) is isomorphic to the structure of the quotient group \(\mathbf{G}' / \mathbf{H}'\) (to see that the condition is sufficient, show first, with the aid of §4, Exercise 2(a), that if \(\top\) is associative it determines on \(\mathbf{G} / \mathbf{H}\) a group structure; denoting by \(\mathbf{K}\) the largest normal subgroup of \(\mathbf{G}'\) contained in \(\mathbf{H}'\) , show then, using the associativity condition on \(\top\) , that \((r(\mathbf{X} \top \mathbf{Y}))^{-1} r(\mathbf{X}) r(\mathbf{Y}) \in \mathbf{K}\) for all \(\mathbf{X}, \mathbf{Y}\) ; deduce that the mapping \(\mathbf{X} \mapsto r(\mathbf{X})\) . \(\mathbf{K}\) is an isomorphism of the group \(\mathbf{G} / \mathbf{H}\) (with the law \(\top\) ) onto the quotient group \(\mathbf{G}' / \mathbf{K}\) ; conclude that \(\mathbf{H}' = \mathbf{K}\) , noting that \(\mathbf{H}'\) is a union of cosets mod. \(\mathbf{K}\) ).
\item
  Conversely, suppose given on a set E an internal law of composition \(\top\) such that every left translation is bijective and there exists \(e \in \mathbf{E}\) such that \(x \top e = x\) for all \(x \in \mathbf{E}\) . Let \(\Gamma\) be the group of permutations of E generated by the left translations \(\gamma_x\) and \(\Delta\) the subgroup of permutations in \(\Gamma\) leaving \(e\) invariant; show that to every left coset X modulo \(\Delta\) there corresponds one and only one element \(x \in \mathbf{E}\) such that \(\gamma_x \in \mathbf{X}\) ; if we write \(r(\mathbf{X}) = \gamma_x\) , the mapping \(x \mapsto \gamma_x\) is an isomorphism of the set E with the law \(\top\) onto the set \(\Gamma / \Delta\) with the law \((\mathbf{X}, \mathbf{Y}) \mapsto r(\mathbf{X})r(\mathbf{Y})\) . \(\Delta\) .
\end{enumerate}

\begin{enumerate}
\def\labelenumi{\arabic{enumi}.}
\setcounter{enumi}{19}
\tightlist
\item
  Let \(G\) be a simple group and \(H\) a subgroup of \(G\) of finite index \(n > 1\) .
\end{enumerate}

\begin{enumerate}
\def\labelenumi{(\alph{enumi})}
\item
  Show that \(\mathbf{G}\) is finite and that its order divides \(n!\) (use Exercise 5).
\item
  Show that \(\mathbf{G}\) is commutative if \(n \leqslant 4\) (use Exercise 10).
\end{enumerate}

\begin{enumerate}
\def\labelenumi{\arabic{enumi}.}
\setcounter{enumi}{20}
\tightlist
\item
  Let \(p(n)\) be the number of conjugacy classes of the symmetric group \(\mathfrak{S}_n\) . Show that \(p(n)\) is equal to the number of families \((x_1, \ldots, x_n)\) of integers \(\geqslant 0\) such that \(\sum_{i=1}^{n} i.x_i = n\) .
\end{enumerate}

* Deduce the identity \(\sum_{n=0}^{\infty} p(n) T^n = \prod_{m=1}^{\infty} \frac{1}{1 - T^m}\) .*

\begin{enumerate}
\def\labelenumi{\arabic{enumi}.}
\setcounter{enumi}{21}
\tightlist
\item
  Let \(A = Z/2Z\) , \(B = S_{3}\) and \(G = A \times B\) . If s is an element of B of order 2, let \(\phi\) be the endomorphism of G whose kernel is equal to B and whose image is \(\{1, s\}\) . Show that the centre of G is not stable under \(\phi\) .
\end{enumerate}

¶ 23. (a) Let s be an element of order 2 in the group \(\mathfrak{S}_n\) which is a product of \(k\) transpositions \(t_1, \ldots, t_k\) with mutually disjoint supports. Show that the centralizer \(\mathrm{C}_s\) of \(s\) in \(\mathfrak{S}_n\) is isomorphic to \(\mathfrak{S}_{n-2k} \times \mathrm{A}\) , where A admits a normal subgroup of order \(2^k\) such that A/B is isomorphic to \(\mathfrak{S}_k\) (observe that every element of \(\mathrm{C}_s\) permutes the fixed points of \(s\) and also the supports of \(t_1, \ldots, t_k\) ).

\begin{enumerate}
\def\labelenumi{(\alph{enumi})}
\setcounter{enumi}{1}
\item
  Show that, if \(n = 5\) or \(n \geqslant 7\) , every element of order 2 in \(\mathfrak{S}_n\) whose centralizer is isomorphic to that of a transposition is a transposition (compare the Jordan-Hölder quotients of these centralizers using Exercise 16). Deduce that every automorphism of \(\mathfrak{S}_n\) permutes the transpositions.
\item
  If \(a, b\) are distinct elements of \([1, n]\) , let \(F_a\) (resp. \(F_{ab}\) ) denote the stabilizer of \(a\) (resp. the fixer of \(\{a, b\}\) ) in \(\mathfrak{S}_n\) and \((ab)\) the transposition \(\tau_{a,b}\) . Let \(a, b, c, d\) be distinct elements of \([1, n]\) . Show that, if \(n \geqslant 5\) , the group generated by \(F_{ab}\) and \(F_{cd}\) is \(\mathfrak{S}_n\) (show that this group contains all the transpositions; for this use the fact that, if \(x \in [1, n] = \{a, b, c, d\}\) , then \((ac) = (cx)(ax)(cx)\) ). Show that the group generated by \(F_{ab}\) and \(F_{ac}\) is \(F_a\) (use the same equation as in the above case).
\item
  Show that an automorphism \(\alpha\) of \(\mathfrak{S}_n\) such that \(\alpha(F_a) = F_a\) for all \(a \in [1, n]\) is the identity (show that \(\alpha\) maps every transposition to itself). Show that, if there exists \(a \in [1, n]\) such that \(\alpha(F_a) = F_a\) , \(\alpha\) is an inner automorphism (show that there exists an inner automorphism \(\beta\) such that \(\beta \circ \alpha\) maps each of the groups \(F_b\) , \(b \in [1, n]\) , to itself).
\item
  Let \(a, b\) be two distinct elements of \([1, n]\) and let \(\mathbf{C}_{ab}\) be the centralizer of \((a, b)\) in \(\mathfrak{S}_n\) . Show that, if \(n \geqslant 5\) , the derived group of \(\mathbf{C}_{ab}\) is of index 4 in \(\mathbf{C}_{ab}\) (use the fact that the derived group of \(\mathfrak{S}_{n-2}\) is \(\mathfrak{A}_{n-2}\) ). Show that \(\mathbf{F}_{ab}\) is the unique subgroup of \(\mathbf{C}_{ab}\) of index 2 which does not contain \((a, b)\) and is not contained in \(\mathfrak{A}_n\) .
\item
  Show that, if \(n \geqslant 5\) , an automorphism \(\alpha\) of \(\mathfrak{S}_n\) which leaves fixed a transposition \((ab)\) is inner. (Show by means of (e) that \(\alpha(\mathrm{F}_{ab}) = \mathrm{F}_{ab}\) ; show equally that, if \(c \notin \{a, b\}\) , \(\alpha(\mathrm{F}_{ac})\) is either of the form \(F_{ax}, x \notin \{a, b\}\) , or of the form \(\mathrm{F}_{bx}, x \notin \{a, b\}\) ; in the first case conclude that \(\alpha(\mathrm{F}_a) = \mathrm{F}_a\) and apply (d); reduce the second case to the first by multiplying \(\alpha\) by the inner automorphism of \(\mathfrak{S}_n\) defined by \((ab)\) .)
\item
  Deduce from (b) and (f) that every automorphism of \(\mathfrak{S}_n\) is inner if \(n = 5\) or \(n \geqslant 7\) .
\end{enumerate}

\begin{enumerate}
\def\labelenumi{\arabic{enumi}.}
\setcounter{enumi}{23}
\tightlist
\item
  \begin{enumerate}
  \def\labelenumii{(\alph{enumii})}
  \tightlist
  \item
    Show that there exists a subgroup H of \(\mathfrak{S}_6\) which is isomorphic to \(\mathfrak{S}_5\) and leaves fixed no element of \([1,6]\) (let \(\mathfrak{S}_5\) operate by inner automorphisms on the set of its subgroups of order 5, which has 6 elements).
  \end{enumerate}
\end{enumerate}

\begin{enumerate}
\def\labelenumi{(\alph{enumi})}
\setcounter{enumi}{1}
\item
  Show that there exists an automorphism \(\sigma\) of \(\mathfrak{S}_{6}\) which maps (12) to (12)(34)(56). (Let \(\mathfrak{S}_{6}\) operate on \(\mathfrak{S}_{6} / \mathrm{H}\) , where \(\mathrm{H}\) is chosen as above.)
\item
  Show that the group of inner automorphisms of \(\mathfrak{S}_{6}\) is of index 2 in the group of all automorphism of \(\mathfrak{S}_{6}\) . (Let \(\alpha\) be an automorphism of \(\mathfrak{S}_{6}\) ; show that \(\alpha\) maps a transposition either to a transposition or to a conjugate of (1 2)(3 4)(5 6). Using Exercise 23, show that in the first (resp. second) case \(\alpha\) (resp. \(\alpha \circ \sigma\) ) is an inner automorphism.)
\end{enumerate}

\begin{enumerate}
\def\labelenumi{\arabic{enumi}.}
\setcounter{enumi}{24}
\tightlist
\item
  Let \(n\) be an integer \(\geqslant 5\) . Show that the subgroups of \(\mathfrak{S}_n\) of order \((n - 1)!\) are isomorphic to \(\mathfrak{S}_{n-1}\) and form a single conjugacy class (resp. two conjugacy classes) if \(n \neq 6\) (resp. if \(n = 6\) ). (Let \(H\) be such a subgroup and \(x_1, \ldots, x_n\) the elements of \(\mathfrak{S}_n / \mathrm{H}\) . Show that the action of \(\mathfrak{S}_n\) on the \(x_i\) defines an automorphism of \(\mathfrak{S}_n\) and apply Exercises 23 and 24.)
\end{enumerate}

¶ 26. (a) Let G be a group operating on two finite sets \(E_{1}\) and \(E_{2}\) . If \(s \in G\) , let \(s_{1}\) (resp. \(s_{2}\) ) denote the permutation of \(E_{1}\) (resp. \(E_{2}\) ) defined by s and \(E_{1}^{s}\) (resp. \(E_{2}^{s}\) ) the set of elements invariant under \(s_{1}\) (resp. \(s_{2}\) ).

Show the equivalence of the following properties:

\begin{enumerate}
\def\labelenumi{(\roman{enumi})}
\item
  For all \(s \in G\) , \(\operatorname{Card}(\mathbf{E}_1^s) = \operatorname{Card}(\mathbf{E}_2^s)\) .
\item
  For all \(s \in \mathbf{G}\) , the orders of the component cycles of \(s_1\) (cf.~no. 7, Proposition 7) are the same (to within a permutation) as those of \(s_2\) .
\item
  For all \(s \in \mathbf{G}\) , there exists a bijection \(f_s: \mathbf{E}_1 \to \mathbf{E}_2\) such that \(s_2 \circ f_s = f_s \circ s_1\) .
\end{enumerate}

If these properties hold, the G-sets \(E_{1}\) and \(E_{2}\) are said to be weakly equivalent.

\begin{enumerate}
\def\labelenumi{(\alph{enumi})}
\setcounter{enumi}{1}
\item
  Give an example of two weakly equivalent G-sets which are not isomorphic (take G to be a non-cyclic group of order 4).
\item
  Let \(H_{1}\) and \(H_{2}\) be two subgroups of a finite group G. Show the equivalence of the following properties:
\end{enumerate}

\begin{enumerate}
\def\labelenumi{(\arabic{enumi})}
\tightlist
\item
  For every conjugacy class C of G,
\end{enumerate}

\[
\operatorname{Card} (\mathrm{C} \cap \mathrm{H} _ {1}) = \operatorname{Card} (\mathrm{C} \cap \mathrm{H} _ {2}).
\]

\begin{enumerate}
\def\labelenumi{(\arabic{enumi})}
\setcounter{enumi}{1}
\tightlist
\item
  The G-sets \(\mathbf{G} / \mathbf{H}_1\) and \(\mathbf{G} / \mathbf{H}_2\) are weakly equivalent.
\end{enumerate}

\begin{enumerate}
\def\labelenumi{(\alph{enumi})}
\setcounter{enumi}{3}
\item
  Show that, if \(H_{1}\) and \(H_{2}\) satisfy properties (1) and (2) above and \(H_{1}\) is normal in G, then \(H_{2} = H_{1}\) .
\item
  Let \(\mathbf{G} = \mathfrak{A}_6\) and let
\end{enumerate}

\[
\begin{array}{l} \mathbf {H} _ {1} = \{e, (1 2) (3 4), (1 3) (2 4), (1 4) (2 3) \} \\ \mathbf {H} _ {2} = \{e, (1 2) (3 4), (1 2) (5 6), (3 4) (5 6) \}. \end{array}
\]

Show that \(H_{1}\) and \(H_{2}\) satisfy properties (1) and (2) and that \(H_{1}\) and \(H_{2}\) are not conjugate in \(A_{6}\) (nor in \(S_{6}\) ).

\begin{enumerate}
\def\labelenumi{\arabic{enumi}.}
\setcounter{enumi}{26}
\item
  Let Y be a subset of a set X and let A be the fixer of Y in the permutation group \(S_{X}\) of X. Let M be the submonoid of \(S_{X}\) consisting of the elements s such that \(sAs^{-1} \subset A\) . Show that M is a subgroup of \(S_{X}\) if and only if one of the sets Y and X = Y is finite.
\item
  Let G be a group, A and B two subgroups of G and \(\phi\) an isomorphism of A onto B. Show that there exists a group \(G_{1}\) containing G such that \(\phi\) is the restriction of an inner automorphism of \(G_{1}\) (use the permutation group of the set G). Show that, if G is finite, \(G_{1}\) can be chosen to be finite.
\item
  Let X be an infinite set. Let G be the subgroup of \(S_{X}\) consisting of the permutations \(\sigma\) which enjoy the following property: there exists a finite subset \(Y_{\sigma}\) of X such that \(\sigma x = x\) if \(x \in X - Y_{\sigma}\) and the restriction of \(\sigma\) to \(Y_{\sigma}\) is even.
\end{enumerate}

\begin{enumerate}
\def\labelenumi{(\alph{enumi})}
\item
  Show that G is a non-commutative simple group. Show that every subset of G which generates G is equipotent to X.
\item
  Show that every finite group is isomorphic to a subgroup of G (use the fact that \(S_{n}\) is isomorphic to a subgroup of \(A_{n+2}\) ).
\end{enumerate}

